\documentclass[11pt,a4paper]{article}

% ---------------------------------------------------------------
% Core mathematics packages
% ---------------------------------------------------------------
\usepackage{amsmath, amssymb, amsthm, mathtools}
\usepackage{bm}            % bold math symbols
\usepackage{dsfont}        % indicator function \mathds{1}

% ---------------------------------------------------------------
% Algorithms
% ---------------------------------------------------------------
\usepackage[ruled,vlined,linesnumbered]{algorithm2e}

% ---------------------------------------------------------------
% Graphics, figures, and plots
% ---------------------------------------------------------------
\usepackage{graphicx}
\usepackage{tikz}
\usepackage{pgfplots}
\pgfplotsset{compat=1.18}
\usepackage{subfigure}     % side-by-side figures
\usepackage{adjustbox}     % fallback rescaling for wide TikZ diagrams

% ---------------------------------------------------------------
% Tables
% ---------------------------------------------------------------
\usepackage{booktabs}
\usepackage{tabularx}
\usepackage{siunitx}       % decimal-aligned columns, SI units

% ---------------------------------------------------------------
% Cross-referencing and citations
% ---------------------------------------------------------------
\usepackage[numbers,sort&compress]{natbib}
\usepackage[colorlinks=true,linkcolor=blue,citecolor=blue,urlcolor=blue]{hyperref}
\usepackage[capitalise,noabbrev]{cleveref}

% ---------------------------------------------------------------
% Typography
% ---------------------------------------------------------------
\usepackage{microtype}
\usepackage{parskip}
\setlength{\parindent}{0pt}
\setlength{\parskip}{6pt}

% ---------------------------------------------------------------
% Page layout
% ---------------------------------------------------------------
\usepackage[margin=2.5cm]{geometry}

% ---------------------------------------------------------------
% Miscellaneous
% ---------------------------------------------------------------
\usepackage{xcolor}
\usepackage{todonotes}     % \todo{} markers for placeholders (paragraph-mode only)
% \todo{} relies on \marginpar and fails in restricted horizontal mode
% (captions, table cells, optional arguments); use \inlinetodo{} there instead.
\newcommand{\inlinetodo}[1]{\textcolor{red}{\textbf{[TODO: #1]}}}

% ---------------------------------------------------------------
% Theorem environments (amsthm)
% ---------------------------------------------------------------
\theoremstyle{plain}
\newtheorem{theorem}{Theorem}[section]
\newtheorem{lemma}[theorem]{Lemma}
\newtheorem{proposition}[theorem]{Proposition}
\newtheorem{corollary}[theorem]{Corollary}

\theoremstyle{definition}
\newtheorem{definition}[theorem]{Definition}
\newtheorem{assumption}{Assumption}
\newtheorem{example}[theorem]{Example}

\theoremstyle{remark}
\newtheorem{remark}[theorem]{Remark}

% ---------------------------------------------------------------
% Notation macros
% ---------------------------------------------------------------
% Norms and inner products
\newcommand{\norm}[1]{\left\lVert #1 \right\rVert}
\newcommand{\normF}[1]{\left\lVert #1 \right\rVert_{\mathrm{F}}}
\newcommand{\ip}[2]{\left\langle #1,\, #2 \right\rangle}
\newcommand{\abs}[1]{\left\lvert #1 \right\rvert}

% Calculus and optimisation
\newcommand{\grad}{\nabla}
\newcommand{\Hess}{\nabla^2}
\newcommand{\Lapl}{\Delta}
\newcommand{\divergence}{\operatorname{div}}
\newcommand{\curl}{\operatorname{curl}}

% Function spaces
\newcommand{\Lp}[1]{L^{#1}}
\newcommand{\Sob}[2]{H^{#1}(#2)}
\newcommand{\SobZ}[2]{H^{#1}_0(#2)}

% Operators and matrices
\newcommand{\bA}{\mathbf{A}}
\newcommand{\bB}{\mathbf{B}}
\newcommand{\bJ}{\mathbf{J}}
\newcommand{\bK}{\mathbf{K}}
\newcommand{\bI}{\mathbf{I}}
\newcommand{\bx}{\mathbf{x}}
\newcommand{\by}{\mathbf{y}}
\newcommand{\bz}{\mathbf{z}}
\newcommand{\bu}{\mathbf{u}}
\newcommand{\bv}{\mathbf{v}}

% Probability and statistics
\newcommand{\E}{\mathbb{E}}
\newcommand{\Prob}{\mathbb{P}}
\newcommand{\Var}{\operatorname{Var}}
\newcommand{\Cov}{\operatorname{Cov}}
\newcommand{\indicator}{\mathbf{1}}

% Asymptotic notation
\newcommand{\bigO}[1]{\mathcal{O}\!\left(#1\right)}
\newcommand{\smallO}[1]{o\!\left(#1\right)}

% Common domains
\newcommand{\R}{\mathbb{R}}
\newcommand{\N}{\mathbb{N}}
\newcommand{\C}{\mathbb{C}}

% Particle-method notation
\newcommand{\xk}{x_k}
\newcommand{\xik}{x_k^{(i)}}
\newcommand{\wik}{w_k^{(i)}}
\newcommand{\Np}{N}
\newcommand{\dt}{\Delta t}

\title{Neural Particle Method}

\author{Osian Shelley\\
  \texttt{osiandshelley@gmail.com}}

\date{\today}

\begin{document}

\maketitle

\begin{abstract}
  \todo{150--250 word abstract. State: the problem being solved (a particle-based
  discretisation of a PDE / measure-valued evolution learned or accelerated with a
  neural network), the method (neural particle method), the theoretical or empirical
  guarantee obtained, and the significance relative to existing particle methods and
  neural PDE solvers.}
\end{abstract}

\section{Introduction}\label{sec:intro}

\todo{Motivate the problem. Particle methods (e.g.\ SPH, vortex methods, interacting
particle approximations of McKean--Vlasov / Fokker--Planck equations) discretise a
measure-valued PDE by a finite collection of particles. State what is added by
combining this discretisation with a neural network component (e.g.\ learned
interaction kernel, learned drift, or learned closure).}

\todo{Survey related work using \texttt{\textbackslash citet\{\}} and
\texttt{\textbackslash citep\{\}}: classical particle methods, mean-field limits,
and neural approaches to PDEs / interacting particle systems.}

\todo{State the gap addressed and list contributions (numbered or itemised).}

\todo{Outline the remainder of the paper.}

\section{Problem Formulation}\label{sec:problem}

\todo{State the PDE or measure-valued evolution of interest, e.g.\ a Fokker--Planck
or McKean--Vlasov equation
\[
  \partial_t \rho_t = -\divergence\bigl(b[\rho_t]\,\rho_t\bigr)
    + \tfrac{1}{2}\Lapl\bigl(\sigma^2 \rho_t\bigr),
\]
together with the empirical particle approximation
\[
  \rho_t \approx \rho_t^\Np \coloneqq \frac{1}{\Np}\sum_{i=1}^{\Np}
    \indicator_{\{X_t^{(i)}\}},
\]
and describe where a neural network parametrises part of the dynamics.}

\section{Neural Particle Method}\label{sec:method}

\todo{Define the particle update rule. Example scaffold:
\[
  \xik[k+1] = \xik[k] + \dt\, b_\theta\bigl(\xik[k], \{\xik[k]\}_{j=1}^{\Np}\bigr)
    + \sigma\sqrt{\dt}\,\xi_k^{(i)},
\]
where $b_\theta$ is a neural network with parameters $\theta$, and describe the loss
used to train $\theta$.}

\begin{algorithm}[H]
\caption{Neural Particle Method}\label{alg:npm}
\KwIn{Initial particles $\{x_0^{(i)}\}_{i=1}^{\Np}$; network $b_\theta$; step $\dt$;
      number of steps $K$}
\KwOut{Particle trajectories $\{x_k^{(i)}\}_{k=0}^{K}$}
\For{$k = 0, 1, \ldots, K-1$}{
  \For{$i = 1, \ldots, \Np$}{
    Evaluate $b_\theta\bigl(\xik[k], \{\xik[k]\}_{j=1}^{\Np}\bigr)$\;
    Update $\xik[k+1] \leftarrow \xik[k] + \dt\, b_\theta(\cdot) + \sigma\sqrt{\dt}\,\xi_k^{(i)}$\;
  }
}
\Return{$\{x_k^{(i)}\}$}\;
\end{algorithm}

\section{Assumptions}\label{sec:assumptions}

\begin{assumption}\label{ass:regularity}
  \todo{State regularity conditions on the true drift / interaction kernel and on
  the network parametrisation $b_\theta$ (e.g.\ Lipschitz continuity, boundedness).}
\end{assumption}

\section{Main Results}\label{sec:results}

\begin{theorem}[Descriptive title, e.g.\ propagation of chaos / convergence rate]\label{thm:main}
% TODO: replace the bracketed title above once the main result is fixed.
  \todo{Under \cref{ass:regularity}, state the main theoretical result: convergence
  of the neural particle approximation to the mean-field limit, or a generalisation
  / approximation error bound.}
\end{theorem}

\begin{proof}
  \todo{Proof to be completed.}
\end{proof}

\section{Numerical Experiments}\label{sec:experiments}

\todo{Describe test problems, baselines (classical particle methods, exact/reference
solutions where available), metrics, and present results as tables (\cref{tab:results})
and figures.}

\begin{table}[ht]
\centering
\caption{\inlinetodo{Comparison against baseline methods.}}\label{tab:results}
\begin{tabular}{lccc}
\toprule
Method & \inlinetodo{Metric 1} & \inlinetodo{Metric 2} & \inlinetodo{Metric 3} \\
\midrule
\inlinetodo{Baseline} & -- & -- & -- \\
\inlinetodo{Neural particle method} & -- & -- & -- \\
\bottomrule
\end{tabular}
\end{table}

\section{Conclusion}\label{sec:conclusion}

\todo{Summarise contributions and results. State limitations and directions for
future work.}

\begin{thebibliography}{99}

\bibitem{golub2013}
G.~H.~Golub and C.~F.~Van~Loan,
\textit{Matrix Computations},
4th~ed.,
Johns Hopkins University Press, Baltimore, MD, 2013.

\bibitem{raissi2019}
M.~Raissi, P.~Perdikaris, and G.~E.~Karniadakis,
``Physics-informed neural networks: A deep learning framework for solving forward
and inverse problems involving nonlinear partial differential equations,''
\textit{Journal of Computational Physics},
vol.~378, pp.~686--707, 2019.

\end{thebibliography}

\end{document}
